
\begin{abstract}
With the capability of modeling bidirectional contexts, denoising autoencoding based pretraining like BERT achieves better performance than pretraining approaches based on autoregressive language modeling.
However, relying on corrupting the input with masks, BERT neglects dependency between the masked positions and suffers from a pretrain-finetune discrepancy.
In light of these pros and cons, we propose XLNet, a generalized autoregressive pretraining method that (1) enables learning bidirectional contexts by maximizing the expected likelihood over all permutations of the factorization order and (2) overcomes the limitations of BERT thanks to its autoregressive formulation.
Furthermore, XLNet integrates ideas from Transformer-XL, the state-of-the-art autoregressive model, into pretraining.
Empirically, under comparable experiment settings, XLNet outperforms BERT on 20 tasks, often by a large margin, including question answering, natural language inference, sentiment analysis, and document ranking.\footnote{Pretrained models and code are available at \url{https://github.com/zihangdai/xlnet}}.

 % (+8.5 on RACE, +6.8 SQuAD v2)
 % (+8.2 on RACE, +2.2 on MNLI)

% Empirically, XLNet consistently outperforms existing methods and achieves state-of-the-art results on 12 tasks including question answering, natural language inference, sentiment analysis, and document ranking, with an improvement over the previous best by 8.2, 2.2, 2.0 points on 
% Notably, XLNet single models beat all the previous ensembles on GLUE, RACE, and SQuAD v1.


% Despite the improvement, BERT suffers from a pretrain-finetune discrepancy due to the corrupted input and is restricted to an independent reconstruction distribution.
% In light of these pros and cons, we propose XLNet, a generalized autoregressive pretraining method that (1) enables models to learn bidirectional contexts by maximizing the expected likelihood over all permutations of the factorization order and (2) naturally avoids the two aforementioned weaknesses of BERT thanks to its auto-regressive formulation.

%Auto-encoding approaches like BERT have been dominant for language representation pretraining. Autoregressive (AR) models have potential advantages over BERT such as not relying on certain independence assumptions, but fall short at modeling bidirectional contexts. We propose XLNet that maximizes the expected likelihood over all permutations of a sequence, which generalizes AR modeling with the ability to learn bidirectional contexts. This unifies AR modeling and pretraining, and enables pretraining to benefit from the AR framework, including leveraging the latest AR progress Transformer-XL. Empirically, XLNet consistently outperforms existing methods and achieves state-of-the-art results on 14 tasks including question answering, natural language inference, sentiment analysis, and document ranking. Notably, XLNet single models are able to beat all the previous ensembles on GLUE, RACE, and SQuAD v1.

% For language model pretraining, autoregressive (AR) models have unique advantages such as weaker assumptions and more rapid community progress. However, AR models lack the ability to learn deep bidirectional contexts and thus underperform auto-encoding methods like BERT. We propose XLNet that generalizes AR modeling by considering all permutations of a sequence, which naturally learns bidirectional contexts and bridges the gap between AR models and pretraining. As a result, XLNet is able to leverage the benefits and progress of AR, including the latest Transformer-XL architecture. Empirically, XLNet consistently outperforms other pretrained models and achieves state-of-the-art results on 14 tasks including question answering, natural language inference, sentiment analysis, and document ranking. Notably, single models of XLNet with a standard parameter budget are able to beat all the previous ensemble results on GLUE, RACE, and SQuAD v1.1.

% Compared to the dominant auto-encoding approach, autoregressive models have unique advantages for language model pretraining such as less simplified assumptions and more rapid community progress. However, there is a gap between existing autoregressive models and pretraining because the ability to learn deep bidirectional contexts is absent. We propose XLNet to bridge such a gap and fully unleash the potential of autoregressive models. By generalizing autoregressive modeling to consider all permutations of a sequence, XLNet naturally learns bidirectional contexts, while enjoying the unique benefits of autoregressive models for pretraining. Among other benefits, we show that Transformer-XL, the latest progress in autoregressive modeling, improves pretraining performance. Empirically, XLNet consistently outperforms other pretrained models and achieves state-of-the-art results on 14 tasks including question answering, natural language inference, sentiment analysis, and document ranking. Notably, single models of XLNet with a standard parameter budget are able to beat all the previous ensemble results on GLUE, RACE, and SQuAD v1.1.

% Autoregressive models target the fundamental problem of density estimation, on which rapid progress has been made. However, as pointed out by BERT, autoregressive models lack the ability to learn deep bidirectional contexts and thus are insufficient for pretraining, which questions their practical usefulness. We propose XLNet to bridge such a gap. By generalizing autoregressive modeling to consider all permutations of a sequence, XLNet naturally learns bidirectional contexts, while enjoying the unique benefits of autoregressive models for pretraining. As an example, we show that the state-of-the-art autoregressive model Transformer-XL improves pretraining performance.


% In the setting of language model pretraining, autoregressive models are not sufficient to learn deep bidirectional dependency, while denoising auto-encoders cause inconsistency between pretraining and finetuning. To overcome these drawbacks, we propose XLNet, a novel pretraining method that generalizes autoregressive modeling by considering all permutation orders of a sequence, which naturally learns bidirectional contexts without adding artificial noise to the inputs.
% Compared to auto-encoders, our method further allows leveraging the recurrence mechanism introduced in Transformer-XL and improves the capability of encoding long text.
% Empirically, XLNet consistently outperforms other pretrained models and achieves state-of-the-art results on 14 tasks including question answering, natural language inference, sentiment analysis, and document ranking. Notably, single models of XLNet with a standard parameter budget are able to beat all the previous ensemble results on GLUE, RACE, and SQuAD v1.1.

\end{abstract}
